\chapter{Circumcision}

\medskip
\noindent\textit{\textbf{Under review.} This chapter is an AI-assisted educational draft; it carries no source citations and is not a substitute for professional medical advice.}

\section*{Definition and Overview}
Circumcision is the surgical removal of the foreskin, the fold of skin that covers the head (glans) of the penis. It is one of the oldest surgical procedures known and is performed for a range of reasons, including religious and cultural tradition, medical necessity, hygiene, and parental or personal choice. In newborn boys it is a short procedure usually done under local anaesthesia; in older children and adults it generally requires a general anaesthetic.

The operation removes the skin covering the glans while preserving the rest of the penile skin, so that the wound heals with a neat scar line. When performed for medical reasons, the aim is to treat a specific problem of the foreskin such as tightness, inflammation, or inability to retract.

\section*{Epidemiology}
Circumcision is one of the most common surgical procedures worldwide. It is nearly universal in some cultures, particularly where it is practised for religious reasons, such as among Jewish and Muslim communities. In the United States it is performed on more than half of newborn boys, while in most of Europe, the United Kingdom, and Australia it is less common and usually done for medical reasons or parental choice.

The medical reasons for circumcision, such as a tight foreskin (phimosis) or recurrent infection, are most often seen in school-age children. Rates of medically indicated circumcision in adults have declined over recent decades, partly because alternative treatments for tight foreskin, such as creams and stretching, are increasingly used first.

\section*{Aetiology and Pathophysiology}
The medical indications for circumcision relate to problems of the foreskin or penis:

\begin{itemize}
    \item \textbf{Phimosis:} a scarred, tight foreskin that cannot be pulled back, causing pain, poor hygiene, or difficulty urinating (ballooning of the foreskin during urination)
    \item \textbf{Recurrent balanitis:} repeated infection or inflammation of the glans and foreskin
    \item \textbf{Balanitis xerotica obliterans (lichen sclerosus):} a chronic skin condition that causes scarring and tightening of the foreskin
    \item \textbf{Paraphimosis:} a retracted foreskin that cannot be returned over the glans, which can obstruct blood flow and is a surgical emergency
    \item \textbf{Recurrent urinary tract infections} in boys with certain urinary abnormalities
    \item \textbf{Religious, cultural, or parental preference} and personal choice are the most common reasons overall
\end{itemize}

A normal, non-scarred tightness in babies and young boys (physiological phimosis) usually resolves on its own and is not a reason for surgery.

\section*{Clinical Features}
In a medical consultation, the features prompting consideration of circumcision relate to the underlying condition rather than the operation itself.

\begin{itemize}
    \item Inability to retract the foreskin (in older boys and men)
    \item Pain or discomfort when urinating, or ballooning of the foreskin
    \item Soreness, redness, itching, or discharge from under the foreskin
    \item Scarring or a whitish, thickened appearance of the tip of the foreskin
    \item A retracted foreskin that becomes stuck, causing a swollen, painful glans
    \item After surgery: mild pain, swelling, and bruising of the penis, settling over several days
\end{itemize}

\section*{Differential Diagnosis}
Conditions of the foreskin and penis that may be confused with each other include:

\begin{itemize}
    \item Physiological phimosis in young boys (normal tightness that resolves spontaneously) vs.\ true pathological phimosis
    \item Balanoposthitis (non-specific inflammation) vs.\ lichen sclerosus vs.\ early penile cancer (rare)
    \item Paraphimosis, which needs urgent treatment
    \item Congenital penile abnormalities such as hypospadias (a urethral opening on the underside of the penis) or a buried penis
    \item Sexually transmitted infections causing discharge or ulceration in sexually active men
\end{itemize}

\section*{When to Seek Medical Care}
See a doctor if the foreskin is persistently tight, painful, or difficult to clean, if there is recurrent soreness or discharge, or if the tip of the penis is scarred or discoloured.

\textbf{Call 000 for:}
\begin{itemize}
    \item A retracted foreskin that cannot be pulled back over the glans (paraphimosis), with a swollen, painful, or blue-tinged glans
    \item Inability to pass urine at all
    \item Heavy bleeding that does not stop after circumcision
    \item Signs of serious infection: spreading redness, severe pain, or fever after surgery
\end{itemize}

\section*{Diagnosis and Investigations}
Circumcision is usually recommended on the basis of history and physical examination alone; special tests are rarely needed.

\begin{itemize}
    \item \textbf{History:} how the foreskin retracts, any pain, discharge, or difficulty urinating, and any history of infection
    \item \textbf{Examination:} inspection of the glans and foreskin, assessment of whether the foreskin can be retracted, and checking the position of the urethral opening
    \item \textbf{Urine test:} if a urinary tract infection is suspected
    \item \textbf{Swabs:} if there is discharge, to identify infection
    \item \textbf{Skin biopsy:} occasionally, if lichen sclerosus or cancer is suspected
    \item \textbf{Pre-operative checks:} general health review and discussion of anaesthetic for older children and adults
\end{itemize}

\section*{Management}
The choice of treatment depends on the reason for circumcision, the age of the person, and the severity of the problem.

\begin{itemize}
    \item \textbf{Observation:} physiological phimosis in young boys usually needs no treatment
    \item \textbf{Medical treatment:} steroid creams and gentle stretching can treat a mildly tight foreskin without surgery
    \item \textbf{Circumcision:} removal of the foreskin using a surgical clamp (e.g.\ Plastibell) or freehand technique, under local anaesthetic in babies or general anaesthetic in older patients
    \item \textbf{Alternative surgery:} preputioplasty (widening the foreskin with a small incision) can be used for mild tightness when the person wants to keep the foreskin
    \item \textbf{Postoperative care:} simple pain relief, keeping the area clean and dry, and avoiding sexual activity until healed in adults
    \item \textbf{Follow-up:} review of wound healing and, in newborns, checking that the result is satisfactory
\end{itemize}

\section*{Complications}
Circumcision is generally very safe, but, like any operation, it carries some risks.

\begin{itemize}
    \item Bleeding, which is usually minor and easily controlled
    \item Infection of the wound, usually mild and treated with cleaning and sometimes antibiotics
    \item Meatal stenosis -- narrowing of the urethral opening, which can occur months later and make urination painful or difficult
    \item Removal of too much or too little skin, causing a tight or bulky result
    \item Adhesions or skin bridges between the glans and the shaft skin
    \item Damage to the glans or urethral opening (rare, and usually avoidable)
    \item A result that is cosmetically not as expected, causing distress
    \item Anaesthetic risks, which are very low for this short procedure
\end{itemize}

\section*{Prognosis and Outcomes}
The outcome of circumcision is excellent in the vast majority of cases. Healing usually takes one to two weeks, and most boys and men have no long-term problems. The procedure effectively resolves the symptoms of phimosis, balanitis, and paraphimosis, and recurrence of these problems is very rare after complete removal of the foreskin. Cosmetic satisfaction is generally high, although it is normal for the scar line to remain visible. As with any operation, the surgeon should explain the expected result and risks before the procedure.

\section*{Prevention and Self-Care}
\begin{itemize}
    \item In babies and young boys, do not force the foreskin back; gentle washing of the outside of the penis is all that is needed
    \item Keep the penis clean and dry, especially in older boys and men with a retractable foreskin
    \item Seek early advice for recurrent soreness or a foreskin that becomes difficult to retract
    \item After circumcision, follow the wound-care instructions given by the surgeon, and attend the follow-up visit
    \item Choose an experienced, qualified surgeon and discuss the risks, benefits, and alternatives before consenting
    \item Parents deciding about newborn circumcision should weigh the small medical benefits against the risks and consider cultural and personal values
\end{itemize}
